%% kapa: exam=1 round=25 year=2014 period=1 form=A subject=civil_law subject_ko=민법 seq=21 number=21
%% source: Q-Net 감정평가사 기출문제 (cst00309) artlSeq=5101207 file=extracted/1008/5101207/2014년 제25회 감정평가사 1차시험 기출문제/1교시(A형).hwp
%% conversion: kordoc 4.24.51 -> tools/gen1.py (auto); review=proofread-sonnet; flags=none
\begin{question}{21}
복대리에 관한 설명으로 옳은 것은? (다툼이 있으면 판례에 의함)
\begin{choices}
  \item 복대리권은 대리인에 의한 수권행위의 철회에 의하여도 소멸한다.
  \item 법정대리인은 본인의 승낙이 있거나 부득이한 사유가 있는 때가 아니면 복대리인을 선임하지 못한다.
  \item 임의대리인이 본인의 지명 없이 복대리인을 선임한 경우에는, 그 부적임을 알고 본인에게 통지나 그 해임을 태만한 때가 아니면 본인에게 대하여 책임이 없다.
  \item 복대리인이 본인을 위한 것임을 표시하지 아니한 때에는, 원칙적으로 대리인을 위한 것으로 본다.
  \item 복대리인이 선임되면 대리인의 대리권은 소멸하고 복대리인만 본인을 대리하게 된다.
\end{choices}
\end{question}
\answer{1}
